% Generated draft: requires mathematical and source review.
\documentclass[11pt]{article}
\usepackage[margin=25mm]{geometry}
\usepackage{fontspec}
\setmainfont{Latin Modern Roman}
\usepackage{amsmath,amssymb,mathtools,unicode-math}
\setmathfont{Latin Modern Math}
\usepackage{enumitem,needspace,xurl,hyperref}
\setlength{\parindent}{0pt}
\setlength{\parskip}{5pt}
\setlength{\emergencystretch}{3em}
\newcommand{\N}{\mathbb{N}}
\newcommand{\Z}{\mathbb{Z}}
\newcommand{\Q}{\mathbb{Q}}
\newcommand{\R}{\mathbb{R}}
\newcommand{\C}{\mathbb{C}}
\newcommand{\F}{\mathbb{F}}
\newcommand{\A}{\mathbb{A}}
\newcommand{\PP}{\mathbb{P}}
\newcommand{\E}{\mathbb{E}}
\newcommand{\eps}{\varepsilon}
\newcommand{\dd}{\,\mathrm{d}}
\begin{document}
\section*{UnsolvedMath ID: 30005477}
\textbf{Source problem code:} catalogue entry
\subsection{Definitions and mathematical statement}
Let \(P=\C[z_1,\ldots,z_d]\) be the polynomial ring appearing in the record. An ideal \(I\subseteq P\) is homogeneous if every homogeneous component of each polynomial in \(I\) also belongs to \(I\).

The supplied question asks whether essential normality of a compressed coordinate tuple associated with \(I\) implies stable division for \(I\). The record does not specify the ambient Hilbert space, its norm, the space onto which the coordinates are compressed, or the precise stable division property. These omissions prevent a unique fully specified operator-theoretic statement.

For clarity, one possible precise formulation can be described conditionally. Suppose \(H\) is a Hilbert space completion of \(P\) in which each coordinate multiplication operator \(M_{z_j}\) is bounded. Write
\[
\mathcal M_I=\overline{I}^{\,H},\qquad Q_I=H\ominus\mathcal M_I.
\]
If the intended compression is to the quotient module, its coordinate operators are
\[
S_j=P_{Q_I}M_{z_j}|_{Q_I},\qquad 1\leq j\leq d.
\]
The tuple \((S_1,\ldots,S_d)\) is essentially normal if
\[
[S_i^*,S_j]=S_i^*S_j-S_jS_i^*\in\mathcal K(Q_I)
\qquad(1\leq i,j\leq d),
\]
where \(\mathcal K(Q_I)\) denotes the compact operators on \(Q_I\).

An exact polynomial stable division condition relative to \(H\) would require a finite homogeneous generating family \(f_1,\ldots,f_r\) for \(I\) and a constant \(C>0\) such that, for every \(f\in I\), there exist \(a_1,\ldots,a_r\in P\) satisfying
\[
f=\sum_{k=1}^{r}a_kf_k,
\qquad
\sum_{k=1}^{r}\|a_kf_k\|_H^2\leq C\|f\|_H^2.
\]
Here the generators and \(C\) are independent of \(f\). This condition is presented as a candidate interpretation; the supplied record does not identify its intended stable division definition.

\textbf{Question, subject to the unresolved specifications above.} For every homogeneous ideal \(I\subseteq P\) in the intended ambient module, does essential normality of its intended compressed coordinate tuple imply its intended stable division property? Under the conditional quotient-module and exact-division interpretations just displayed, this asks whether the compactness conditions imply the displayed uniform division estimate. The choice of \(H\), the compression convention, and the division convention must be recovered from the cited source before this conditional formulation can be identified with the source problem.
\subsection{Short English statement}
The question asks whether essential normality of the coordinate operators associated with a homogeneous polynomial ideal forces uniformly controlled polynomial division by a fixed finite generating family. The supplied record leaves the Hilbert space, compression convention, and stable division definition unspecified, so the displayed formulation is conditional and requires human mathematical review.
\subsection{Sources}
\begin{itemize}
\item Supplied local UnsolvedMath registry record, permanent ID \(30005477\), title: \emph{Stable Division and Arveson--Douglas Essential Normality}; exact source problem code: \texttt{catalogue entry}. The record assigns the year 2023 and names no proposer.
\item \emph{New Directions in Real Algebraic Geometry} (2023), Oberwolfach Reports, pp. 813--881, DOI \texttt{10.4171/owr/2023/15}: \url{https://doi.org/10.4171/owr/2023/15}. The supplied background describes these reports as documenting mathematical workshops and their open-problem discussions. This reference is also supplied under the literature-source label ``OWR report.''
\item Shalit, \emph{Stable polynomial division and essential normality}: \url{https://doi.org/10.1112/jlms/jdq054}. This is a supplied literature reference; no additional bibliographic details are provided.
\item Literature assessment dated exactly \texttt{2026-08-22}: the supplied curation labels the question open, reports stable division as sufficient for essential normality, reports many known cases of essential normality, and describes the converse in this record as remaining open. These are supplied dated curation claims, not independently verified assessments. The background expressly requests independent verification of time-sensitive status and cited claims.
\item The registry labels the statement \texttt{corrected\_verified} and states that it was reconstructed and checked against the official Oberwolfach report. This is supplied verification metadata, not verification performed for this draft. The retained original passage refers to ``this question'' as the first part of the Arveson--Douglas essential normality conjecture and says that a converse is open; its antecedent is absent from the supplied passage.
\end{itemize}
\end{document}
